% Exemplo: esqueleto de TCC com a classe abntex2 e referências com biblatex-abnt
% Compilar: latexmk -pdf tcc-esqueleto.tex         Testado com: abntex2, biblatex-abnt, TeX Live 2024+
% Mostra a ordem dos elementos (NBR 14724): capa, folha de rosto, resumos, listas, sumário,
% texto, referências, apêndices e anexos. A folha de aprovação assinada e a ficha catalográfica
% costumam ser incluídas como PDF (\includepdf, pacote pdfpages).
%%BEGIN:tccclasse
\documentclass[12pt, a4paper, oneside, english, brazilian,
  sumario=tradicional]{abntex2}
%%END:tccclasse
%%BEGIN:tccpacotes
\usepackage[T1]{fontenc}
\usepackage{lmodern}
\usepackage{indentfirst}             % recua o primeiro parágrafo da seção
\usepackage{graphicx}
\usepackage{csquotes}
\usepackage[backend=biber, style=abnt, uniquename=init, giveninits,
            repeatfields]{biblatex}
\usepackage{ragged2e}
\usepackage{xurl}
\addbibresource{../referencias/catalogo.bib}
\appto{\bibsetup}{\RaggedRight}      % à esquerda (NBR 6023), com hifenização
\urlstyle{same}
\hypersetup{hidelinks}               % links sem moldura (hyperref da abntex2)
%%END:tccpacotes

%%BEGIN:tccdados
\titulo{Meu primeiro TCC com abn\TeX2}
\autor{Nome do Autor}
\local{Cidade}
\instituicao{Nome da Instituição \par Nome do \emph{campus} ou da unidade}
\data{\the\year}
\orientador{Prof. Dr. Fulano de Tal}
\coorientador{Profa. Dra. Beltrana da Silva}
\tipotrabalho{Trabalho de Conclusão de Curso (Graduação)}
\preambulo{Trabalho apresentado à Nome da Instituição como parte das
exigências para a obtenção do título de Bacharel em Nome do Curso.}
%%END:tccdados

\begin{document}

%%BEGIN:tccpretextuais
\pretextual
\imprimircapa
\imprimirfolhaderosto

\begin{resumo}
Texto do resumo em português, em parágrafo único.

\textbf{Palavras-chave}: LaTeX. abnTeX2. Trabalho acadêmico.
\end{resumo}

\begin{resumo}[Abstract]
\begin{otherlanguage*}{english}
Abstract text in English.

\textbf{Keywords}: LaTeX. abnTeX2. Academic writing.
\end{otherlanguage*}
\end{resumo}

\listoffigures*       % o * deixa a lista fora do sumário
\cleardoublepage      % cada lista e o sumário em página própria
\listoftables*
\cleardoublepage
\tableofcontents*
\cleardoublepage
%%END:tccpretextuais

%%BEGIN:tcctextual
\textual
\chapter{Introdução}\label{cap:introducao}

Os objetivos do trabalho são:
\begin{alineas}
  \item levantar os requisitos;
  \item desenvolver o protótipo;
  \item avaliar o protótipo com usuários.
\end{alineas}

Segundo \textcite{ex-livro}, a citação direta longa (mais de três linhas)
fica recuada:
\begin{citacao}
Texto citado literalmente, com mais de três linhas, que a classe abntex2
formata com recuo de 4\,cm, fonte menor e espaçamento simples, como pede a
NBR 10520. A referência vem ao final da citação, entre parênteses
\cite[p.~35]{ex-livro}.
\end{citacao}
%%END:tcctextual

%%BEGIN:tccpostextual
\postextual
\printbibliography[title=Referências]

\begin{apendicesenv}
\chapter{Questionário aplicado}
Texto do apêndice (elaborado pelo autor).
\end{apendicesenv}

\begin{anexosenv}
\chapter{Regulamento do curso}
Texto do anexo (elaborado por terceiros).
\end{anexosenv}
%%END:tccpostextual

\end{document}
